\begin{figure}[htbp]
\centering
\begin{tikzpicture}[font=\small,>=Stealth]

\fill[blue!7] (0,0) rectangle (4,2.4);
\draw[dashed,thick] (0,0) rectangle (4,2.4);
\filldraw[fill=blue!25,draw=blue] (2,1.2) ellipse (1.05 and .65);
\node at (2,1.2) {$\operatorname{supp}\eta$};
\node at (2,3) {interior chart};
\node[align=center] at (2,-.8) {zero near every artificial face\\local boundary integral $=0$};
\fill[blue!7] (6,0) rectangle (10,2.4);
\draw[dashed,thick] (6,0)--(6,2.4)--(10,2.4)--(10,0);
\draw[very thick,red!70!black] (6,0)--(10,0);
\filldraw[fill=blue!25,draw=blue] (7,0)--(9,0) .. controls (9,1.8) and (7,1.8) .. (7,0);
\node at (8,.55) {$\operatorname{supp}\eta$};
\node at (8,3) {boundary chart};
\node[align=center] at (8,-.8) {genuine face $x_n=0$ remains\\local boundary integral survives};
\node at (5,-1.8) {$\eta=\rho_\alpha\omega:\quad\text{artificial faces vanish; genuine faces reassemble into }\partial M$};

\end{tikzpicture}
\caption{Local Stokes applied to a chart-supported form. Compact support removes artificial boundary terms. Smooth extension across the genuine face need not be zero, and this is exactly where the manifold boundary integral comes from.}
\label{fig:local-global-stokes}
\end{figure}
